For completeness, we describe the NTK baseline we implement from
\citet{nichani2024understandingfactualrecalltransformers}.

Let $\mathbf{K}\in\mathbb{R}^{F\times d}$ be the matrix whose $i$th row is $\mathbf{k}_{i}^\top$, and let
$\mathbf{C}_{f}\in\mathbb{R}^{F\times d_c}$ be the matrix whose $i$th row is $\mathbf{c}_{f(i)}^\top$,
where typically $\mathbf{c}_{j}=\mathbf{v}_{j}$. Given hidden width $m$, sample random gate directions
$\mathbf{w}_{r}\stackrel{i.i.d.}{\sim}\mathcal{N}(0,I_{d})$ and random output directions
$\mathbf{p}_{r}\in\mathbb{R}^{d_c}$ with $\|\mathbf{p}_{r}\|_{2}=1$ for $r=1,\dots,m$. Writing
$\mathbf{W}_{\mathrm{gate}}\in\mathbb{R}^{m\times d}$ for the matrix with rows $\mathbf{w}_{r}^\top$ and
$\mathbf{P}=[\mathbf{p}_{1},\dots,\mathbf{p}_{m}]\in\mathbb{R}^{d_{c}\times m}$, define the degree-$1$ Hermite feature
matrix
\[
\mathbf{H} \;:=\; \mathrm{He}_{1}(\mathbf{K}\mathbf{W}_{\mathrm{gate}}^\top)
\;=\;
\mathbf{K}\mathbf{W}_{\mathrm{gate}}^\top \in \mathbb{R}^{F\times m},
\]
where $\mathrm{He}_{1}(t)=t$ is applied entrywise. The up-projection is then chosen as
\[
\mathbf{W}_{\mathrm{up}}
\;:=\;
\frac{1}{m}\big(\mathbf{H}\odot(\mathbf{C}_f\mathbf{P})\big)^\top \mathbf{K} \in \mathbb{R}^{m\times d}.
\]

The NTK MLP construction is then
\[
\hat{\mathbf{c}}_{\mathrm{NTK}}(\mathbf{x})
\;=\;
\mathbf{P}\Big(\sigma(\mathbf{W}_{\mathrm{gate}}\mathbf{x})\odot(\mathbf{W}_{\mathrm{up}}\mathbf{x})\Big).
\]
Throughout this work, we take $\sigma=\mathrm{ReLU}$.

Note that \citet{nichani2024understandingfactualrecalltransformers}'s construction requires choosing a Hermite degree $k$.
This choice plays a role analogous to our kernel choice: it determines which degree polynomial interactions are emphasized by the construction.
In the most favorable case $k=1$, the fit uses the linear features $\mathrm{He}_{1}(\mathbf{K}\mathbf{W}_{\mathrm{gate}}^\top)=\mathbf{K}\mathbf{W}_{\mathrm{gate}}^\top$, and the realized finite-width model remains a gated bilinear MLP.
Thus, in expectation over the random features, the $k=1$ baseline captures quadratic interactions, like how our sketched bilinear construction approximates a quadratic kernel.
This result, and the fact that the capacity bounds in~\citet{nichani2024understandingfactualrecalltransformers} degrade exponentially with $k$, makes using $k=1$ the fairest comparison to our sketched-$K_{2}$ construction.
The main difference is that the NTK construction folds the target codes into the hidden coefficients through $\mathbf{C}_f\mathbf{P}$, whereas our construction uses an explicit bilinear random-feature map followed by a Hebbian readout.
